\chapter{质：属性编码、活动分布与任务读出}
\label{chp:qualia_activation}

关系结构让系统知道对象如何组织，属性内容让系统知道这些位置上承载着什么。我们现在转向“质”的计算表达。面对同一个杯子，系统可以记录颜色、温度、容量、用途与当前可信度；面对同一个事件，系统可以记录参与者、时刻、来源与结果。这些内容并不具有相同类型，也不以相同速度变化。我们需要把它们组织成能够参与计算的属性状态，同时保留各项内容的解释与来源。

我们沿用 Qualia 这个名称指称模型中的属性内容，但不由一个属性向量推断主观感受已经被实现。红色的编码、关于红色的分类结果与红色体验的解释属于不同论证层级。本章研究前两类对象怎样进入状态动力学，体验问题在其专门讨论中另行提出条件。我们的目标是说明内容如何编码、如何被激活、如何在关系上保持身份，以及如何通过任务读出发挥作用。

\section{属性的类型、尺度与来源}
\label{sec:qualia_definition}

我们首先把属性看作带有类型的记录，而不是一个可以任意相加的数值集合。一个位置上的内容可以写为 $x_i=(x_i^{(1)},\ldots,x_i^{(p)})$，其中每一分量具有给定的属性空间。类别、实数测量、文本描述、置信分布和时间记录分别进入对应的空间。我们只有在给出编码方式后，才把它们组合为用于计算的向量。组合向量的维度一致性不意味着原始属性具有同一含义，也不意味着所有分量都可以使用同一种距离来比较。

以杯子为例，温度读数可以带有摄氏度单位，颜色可以是类别或光谱特征，容量可以是带误差范围的数值估计。若把温度数值与颜色类别的内部编号直接求和，计算虽然可以执行，结果却没有明确的属性解释。我们因此为各类属性设置编码函数 $e_a:\mathcal X_a\to\mathbb{R}^{d_a}$，再通过拼接或声明了意义的融合函数得到计算特征。需要返回测量值时，系统保留原始记录或使用经过验证的解码接口，而不是仅凭隐藏向量反推出精确的物理量。

尺度选择会影响模型行为。对一个连续属性采用标准化 $\widetilde x=(x-\mu)/\sigma$ 时，我们要求 $\sigma>0$，并记录参数来自哪些数据与使用条件。同样的数值在不同标准化参数下对应不同的原始量。若系统换用新的参数，却继续读取旧编码而没有迁移标记，历史属性就可能被误解释。标准化参数、编码版本与原始单位因此属于属性记录的辅助信息，它们帮助不同模块在交换内容时保持一致。

属性来源同样重要。一个温度来自传感测量，另一个来自文本推断，还有一个来自旧记录的预测，三者都可以成为当前估计的输入，但不应拥有相同的证据身份。我们可以为记录保存来源键、观测时刻、推断方式与有效区间。系统在发生冲突时，能够据此组织进一步核查。若来源被压缩掉，即使向量保留了一个看似合理的数值，系统也可能无法解释为什么相信它，以及这个信念应当在什么条件下更新。

缺失属性也不是零值。容量未知与容量为零具有不同任务后果，温度没有观测与温度处于零摄氏度也不能共用一个无说明的数值。我们可以采用有效性掩码、未知类别或显式的缺失模型，使计算模块能够区分“没有内容”与“内容恰好为零”。对于概率模型，缺失值的处理还需要与条件分布相匹配；对于确定性规则，系统需要说明缺失是否阻止行动，或者是否触发补充观测。

属性可以随任务重新组织，但重新组织需要保持对应。系统可能在粗粒度任务中只记录“热或不热”，在精细控制中保留完整温度估计。由精细记录得到粗类别通常可以通过阈值读出，而由粗类别恢复精确数值通常没有唯一答案。我们把这种信息方向差异写入接口，避免把压缩后的属性当成未经损失的原始记录。属性表达的有效性由任务需要决定，压缩与解码都必须说明保留了哪些差异。

通过这些约定，“质”成为一种可以维护的计算内容。它不仅包括用于模型运行的特征，也包括理解特征所需的类型、单位、来源与有效性条件。我们不要求每一个内部活动都携带全部辅助记录，而是要求在压缩、传递与读出时存在能够恢复必要解释的路径。这样，活动状态可以保持计算效率，系统又不至于在跨模块交换中丢失内容的身份与可信依据。

\section{编码空间与概率空间的区分}
\label{sec:qualia_elements}

我们用向量表示属性，是为了支持连续计算，但一个向量并不自动是概率分布。设 $z\in\mathbb{R}^d$ 是特征编码，它可以包含正值、负值以及不同尺度的分量。除非我们给出非负性、归一化与事件解释，否则不能把这些分量直接用于信息熵或概率采样。特征向量描述计算坐标，概率分布描述指定事件的不确定性。它们可以通过读出相连接，但属于不同类型的对象。

对于一个具有 $K$ 个候选类别的属性，我们可以先得到分数 $a\in\mathbb{R}^K$，再定义
\[
p_j=\frac{\exp(a_j/\tau)}{\sum_{\ell=1}^K\exp(a_\ell/\tau)},\qquad \tau>0.
\]
由指数函数为正可知 $p_j>0$，由分母定义可知 $\sum_jp_j=1$。这使 $p$ 成为候选类别集合上的概率向量。为了数值稳定，我们可以同时减去 $\max_ja_j$；分子和分母获得相同因子，理论上的分布不变。这里的尺度参数调节分数差异的影响，不需要赋予物理温度单位。

概率向量的归一化并不保证它与实际频率一致。如果模型在一些输入上给某类别较高概率，却经常判断错误，我们还需要检查它在相应数据上的校准与区分能力。归一化只保证概率表达的基本形式，经验可信度需要通过观测建立。我们因此区分编码是否类型正确、分布是否归一化、预测是否准确以及不确定性是否与错误风险一致。任何一个层面通过检查，都不能代替另外几个层面的检查。

在需要比较属性时，我们也要说明距离的作用。连续特征可以使用欧氏距离或学习得到的相似函数，但相似不等于对象相同。两个杯子可以具有接近的颜色编码，却具有不同容量；同一个杯子在光照改变后，视觉编码也可能移动。若系统根据相似度维护身份，就需要加入时间、位置或其他证据。向量空间支持候选检索，身份判断则需要指定决策条件，这种职责划分有助于防止近邻检索变成未经核查的事实合并。

信息熵只能在给定分布上定义。对上述类别分布，我们可写 $H(p)=-\sum_jp_j\log p_j$，并采用 $0\log0=0$ 的延拓约定。它测量分布在这些类别之间的分散程度，而不是直接测量系统的全部认知复杂度。一个分布熵较低，可能表示模型集中相信某个类别，也可能表示它过早排除了正确候选。若要把熵用于控制，我们需要说明它与任务风险之间的关系，并配合误差、来源与预算等观测。

编码空间还可以采用不同的组织形式。一个属性可以由单个向量承载，也可以由多个局部位置共同承载；多个属性可以分通道保存，也可以由一个联合编码表达。联合编码有利于处理属性间的相关性，但可能降低单项编辑的可解释性。我们可以通过属性交换、局部遮挡与重建测试检查这种代价。一个能够重建整体输入的编码，未必能够让任意属性被独立编辑，因此不同的使用目的需要不同的验证。

这些区分使我们能够在后续章节中准确使用范数、距离、概率与熵。范数描述给定编码中的活动大小，距离描述指定空间中的差异，概率描述候选事件的不确定性，熵描述该分布的分散程度。我们不会因为它们都能被计算成一个实数，就把它们合并为同一种“能量”。只有为具体机制建立了明确转换与量纲约定之后，系统才可以把这些观测共同用于状态调节。

\section{激活机制、选择权重与持续更新}
\label{sec:activation_field}

属性记录可以长期存在，但不是所有内容都在每一时刻参与同样强度的计算。我们用激活变量描述当前可用或被选择的程度，用内容变量描述该位置承载的属性。令 $r_i$ 表示被保存的内容，$g_i\in[0,1]$ 表示当前门控权重，可以定义工作表示 $x_i=g_i r_i$。这里的乘法说明门控如何影响内容进入当前计算，而不是宣布内容本身的事实值发生了同比变化。温度记录被降低注意权重，不意味着杯子的实际温度下降。

在这个模型中，我们保留 $r_i$ 与 $g_i$ 的区分，因为仅观察乘积 $x_i$ 通常无法唯一恢复它们。对于非零内容，较大的内容幅度与较小的门控可以产生同样乘积；对于零门控，内容可以继续存在却无法从当前活动中被直接读出。若工程系统需要解释某项信息为何没有影响行动，就要观察门控、来源和访问过程，而不能只根据工作向量中该信息不明显，就断言它已经从全部记忆中消失。

一种简单的门控更新为 $\dot g_i=-\lambda_i(g_i-b_i)$，其中 $\lambda_i>0$、$b_i\in[0,1]$ 在所考察时间区间内固定。解为 $g_i(t)=b_i+(g_i(0)-b_i)e^{-\lambda_it}$。当初值也位于单位区间时，解是初值与目标值的凸组合，因此始终位于单位区间，并逐渐接近 $b_i$。这个结论依赖固定目标；当控制不断改变 $b_i$ 时，我们需要把目标变化纳入分析。门控有界只说明权重合法，并不保证选择的内容正确。

另一种选择机制是在候选内容间分配归一化权重。对于分数 $a_i$，我们采用前一节的 softmax 定义。其偏导满足 $\partial p_i/\partial a_j=\tau^{-1}p_i(\delta_{ij}-p_j)$，由商法则即可得到。这个表达说明，提高一个候选的分数不仅改变它自身的权重，也通过归一化影响其他候选。我们可以据此分析分数误差如何影响选择，但不能把每个权重的改变解释为彼此独立的局部过程。

归一化选择适合分配某一种有限选择权，但它不一定适合表示全部内容活动。两个相互独立且都重要的证据，可能需要同时保持较高激活；把它们强制放入同一个总和为一的集合，会引入不必要的竞争。我们可以把独立门控与局部归一化组合使用，并声明哪些对象共享预算。系统中的竞争来自具体约束，而不是来自所有内容必然相互排斥的假定。这种区分也为 CCM 的工作状态管理提供了更准确的观测对象。

持续更新需要处理输入、衰减与反馈。新证据可以提高某项内容的访问优先级，时间间隔可以改变它的当前相关性，任务控制可以抑制无关内容的进入。我们可以把这些因素写入门控目标或活动更新函数，同时保存控制依据的状态版本。若一个反馈使用旧分布来调节新分布，就需要考虑延迟。这里的控制作用是改变计算中内容的可用程度及其相互影响，不必把它描述为外加的物理力。

在文档任务中，关键证据可以保留在外部记录中，并在工作状态中只保持来源键、摘要与当前选择权重。当摘要不足以回答问题时，系统重新访问原始证据。激活下降因此可以释放当前工作资源，而不必等同于删除全部历史。AgentOS 可以通过边界与卸载机制组织这种交换；HSF-HD 需要分析活动变化、访问延迟与内容保持之间的耦合。我们以这些可观察机制解释选择与持续性，而不从一个活动峰值推断认知体验已经被完整说明。

\section{内容读出、属性干预与记忆接口}
\label{sec:qualia_function}

内容通过任务读出发挥作用。令当前联合表示为 $(S,X)$，输出为 $\widehat y=g(S,X)$。对于杯子任务，输出可以是对象位置、可用容量估计或一个带前提的动作建议；对于交付事件，输出可以是物品当前持有者与支持该判断的来源。读出不是把全部隐藏状态逐项翻译成文字，而是根据任务要求提取必要结果。我们需要说明哪些差异应当影响输出，以及哪些仅由内部编号或编码尺度引起的差异不应改变结论。

一种有效检查是对属性实施干预。在固定关系结构下改变杯子的颜色，容量判断不应因为颜色编码的任意变化而无依据地改变；在固定对象身份下改变容量估计，关于盛水数量的预测则应有相应变化。我们可以定义属性编辑 $T_a$，比较 $g(S,T_aX)$ 与任务所要求的结果。编辑后的结果需要符合被改变属性的作用范围。若一个颜色操作同时改变多个不相关对象的温度记录，问题可能来自属性绑定、联合编码或编辑接口。

干预还要区分事实修改与假设推演。我们可以在一个暂存状态中把杯壁设为破损，分析可能后果，而不立即把真实杯子的记录改成已破损。系统通过分支身份与来源标记区分当前事实、预测状态与反事实条件。这样，推演结果可以参与计划，却不会污染已有证据。对于长期运行的智能体，这种区分尤其重要，因为一个未经执行的计划可能在多次总结之后被误写成已经发生的经历。

内容的时间更新需要保持事件顺序。一个对象的当前属性来自最近有效观测或经过声明的估计机制；情景记录保存观测与变化发生的过程；持久结构保存可复用的编码方式、关系规则与技能。我们可以把这些作用分别映射到 AgentOS 的工作记忆 $h(t)$、情景记忆 $E$ 和结构记忆 $\Theta$。这是一种工程实例化，不意味着 HSF-HD 的一般内容变量只允许三种存储形态，也不要求所有实现把三个对象放入相同位置。

当内容来自外部记忆 $M_e$ 时，我们还需要处理访问代价与有效性。外部记录可以保存更完整的证据，但调用具有延迟，返回结果可能过时，访问也可能失败。工作状态可以先保存摘要与引用，在需要精确信息时再召回。这个设计改变了当前活动与可访问内容之间的关系：某项内容没有完整进入工作状态，不代表系统无法获得它；它能被获得，也不代表能够在当前预算与时限内使用。边界管理需要把这些条件纳入调度。

我们因此分别观察保存质量、召回质量、读出质量和行动后果。保存质量关注属性身份、来源与时序是否被维护；召回质量关注当前任务是否找到必要内容；读出质量关注内容是否被正确使用；行动后果则检验系统是否在约束内实现目标。一个任务失败可能发生在任何环节。若仅用最终输出评价全部内容机制，我们就难以判断是内容丢失、访问错误、绑定错误，还是控制选择没有把正确内容送入当前计算。

本章的“质”由此具有完整的计算含义：它是带有类型与来源的属性内容，可以被编码、选择、传播、保存与读取。激活变量决定当前怎样使用内容，结构变量决定内容怎样组织，记忆接口决定过去与外部内容怎样进入当前过程。我们通过明确的运算与验证讨论这些作用，不把数值强度直接等同于物理能量，也不把内容向量直接等同于主观体验。接下来，第一卷将进一步讨论联合状态空间及绑定形式，使关系与内容能够在同一个模型中持续演化。
